\begin{table}[!tbp]\centering
\caption{Primera etapa: migración del pico y cumplimiento del salario mínimo}\label{tab:firststage}
\begin{threeparttable}
\footnotesize
\begin{tabular}{lccc}
\toprule
  & Ventana previa & Ventana posterior & \\
\midrule
\multicolumn{4}{l}{\emph{Panel A: proporción de asalariados formales a $\pm$2.5\% de cada mínimo}} \\
Todos los asalariados formales & & & \\
\quad En el mínimo anterior (930) & 0.114 & 0.011 & \\
\quad En el nuevo mínimo (1025) & 0.052 & 0.117 & \\
Asalariados formales de baja educación & & & \\
\quad En el mínimo anterior (930) & 0.142 & 0.016 & \\
\quad En el nuevo mínimo (1025) & 0.064 & 0.140 & \\
\midrule
\multicolumn{4}{l}{\emph{Panel B: DiD sobre la proporción que gana menos de 1025 (umbral fijo)}} \\
 & DiD & (s.e.) & $p$ tend. previa \\ \midrule
Asalariados formales & -0.050 & (0.030) & 0.099 \\
Asalariados informales & 0.042$^{***}$ & (0.012) & 0.068 \\
Todos los asalariados & 0.025 & (0.021) & 0.007 \\
\bottomrule
\end{tabular}
\begin{tablenotes}\footnotesize
\item Notas: El panel A reporta la proporción ponderada de asalariados privados formales (remuneración mensual equivalente) a $\pm$2.5 por ciento de cada mínimo legal, en ventanas simétricas de tres trimestres antes (2021Q3--2022Q1) y después (2022Q3--2023Q1) de la reforma. El panel B reporta estimaciones DiD de $\mathrm{Low}\times\mathrm{Post}$ para un indicador de remuneración mensual equivalente menor a 1025 soles (umbral fijo en todos los períodos), según condición de formalidad, con errores estándar agrupados por departamento y el valor $p$ conjunto de tendencias previas en la última columna. Factores de expansión de la ENAHO en todos los casos. $^{*}p<0.1$, $^{**}p<0.05$, $^{***}p<0.01$.
\end{tablenotes}
\end{threeparttable}\end{table}
